%=====================================================================
%  IEEE journal manuscript. Compile with pdfLaTeX + BibTeX
%  Order:  pdfLaTeX -> BibTeX -> pdfLaTeX -> pdfLaTeX
%=====================================================================
\documentclass[journal]{IEEEtran}
\IEEEoverridecommandlockouts
\usepackage{cite}
\usepackage{amsmath,amssymb,amsfonts}
\usepackage{algorithmic}
\usepackage{graphicx}
\usepackage{textcomp}
\usepackage{xcolor}
\usepackage{booktabs}
\usepackage{array}
\usepackage{multirow}
\usepackage{url}
\def\BibTeX{{\rm B\kern-.05em{\sc i\kern-.025em b}\kern-.08em T\kern-.1667em\lower.7ex\hbox{E}\kern-.125emX}}

% HAG-Net headline numbers (single source of truth)
% Under the residual-fallback rule, the learnable correction is rejected on
% validation for every fold on this cohort (gate -> small negative, val mF1 does
% not exceed the prior), so HAG-Net's staging output equals the classical prior,
% whose rigorous ten-fold numbers (results/ensemble7_v2_all.json) are reported.
\newcommand{\hagAccR}{0.738}   % no-HMM accuracy   (prior no-HMM: 0.7375)
\newcommand{\hagMfR}{0.680}    % no-HMM macro-F1   (prior no-HMM: 0.6803)
\newcommand{\hagKR}{0.634}     % no-HMM kappa      (prior no-HMM: 0.6342)
\newcommand{\hagAccH}{0.746}   % +HMM accuracy     (prior +HMM: 0.7464)
\newcommand{\hagMfH}{0.675}    % +HMM macro-F1     (prior +HMM: 0.6753)
\newcommand{\hagKH}{0.641}     % +HMM kappa        (prior +HMM: 0.6415)

\begin{document}

\title{HAG-Net: A Hemispheric Asymmetry-Guided Hybrid Network for Interpretable Sleep Staging and Lesion-Severity Biomarkers in Subacute Ischemic Stroke}

\author{Md Wahiduzzaman Sua, Esm-e Moula Chowdhury Abha, and Md Imtiaj Alam Sajin%
\thanks{The authors are with the Department of Computer Science and Engineering,
American International University-Bangladesh, Dhaka, Bangladesh
(e-mail: 26-94088-2@student.aiub.edu; 26-94089-2@student.aiub.edu;
26-94090-2@student.aiub.edu).}}

\markboth{Manuscript, July 2026}%
{Sua \MakeLowercase{\textit{et al.}}: HAG-Net for Interpretable Sleep Staging in Ischemic Stroke}

\maketitle

%=====================================================================
\begin{abstract}
Automated sleep staging is mature on healthy cohorts but degrades sharply in
focal brain injury, where lesions distort the very micro-events (spindles, slow
waves, ocular movements, muscle atonia) on which scoring depends. On iSLEEPS,
the first public polysomnography corpus of subacute ischemic stroke (99 patients,
95{,}305 epochs), scaling model capacity fails: every convolutional, recurrent,
transfer-learned, and raw-signal deep network we test underperforms classical
features, the signature of a data-limited rather than capacity-limited regime.
We therefore propose \textbf{HAG-Net}, a hemispheric asymmetry-guided hybrid that
is deep where deep helps and classical where classical helps. HAG-Net couples
(i) a physiologically grounded dual feature representation, (ii) a strong
gradient-boosting \emph{classical prior} stream, (iii) a novel
\emph{hemispheric-asymmetry graph} module that attends over a montage graph and
pools signed differences of homologous derivations to expose the lesion signal,
(iv) a bidirectional selective state-space temporal decoder, and (v) a
residual-gated fusion anchored at the classical prior by construction, with
hidden-Markov grammar decoding. Under the dataset's patient-exclusive ten-fold
protocol, HAG-Net matches the published deep state of the art on accuracy
($\hagAccH$ vs.\ $0.747$) and \emph{exceeds} it on the minority-sensitive metrics
(Cohen's $\kappa$ $\hagKH$ vs.\ $0.640$; macro-F1 up to $\hagMfR$ vs.\ $0.677$ at
the no-smoothing operating point), while subject-paired significance tests confirm
that it dominates every pure-deep baseline ($p<10^{-9}$). On this 99-patient cohort
the deep streams do not raise staging accuracy; the residual gate holds performance
at the strong classical prior, and their value lies in the asymmetry biomarker and
the calibrated uncertainty they add.
Because HAG-Net is interpretable by construction, its
asymmetry stream doubles as a clinical instrument: inter-hemispheric spindle
asymmetry scales with stroke severity (Spearman $\rho=0.41$, $p=0.006$), and a
split-conformal layer yields calibrated prediction sets (empirical coverage
$0.900$ at $\alpha{=}0.1$; expected calibration error $0.038$). Code and features
are released.
\end{abstract}

\begin{IEEEkeywords}
Sleep staging, electroencephalography, ischemic stroke, polysomnography, graph
attention, state-space models, gradient boosting, conformal prediction,
interpretable machine learning, biomarker.
\end{IEEEkeywords}

%=====================================================================
\section{Introduction}
\IEEEPARstart{S}{leep} is scored in 30-second epochs into five stages, Wake (W),
N1, N2, N3 and rapid-eye-movement (REM), following the American Academy of Sleep
Medicine (AASM) rules~\cite{aasm2023manual}. Manual scoring is slow, costly, and
bounded by inter-rater disagreement, which has driven two decades of automation.
Modern deep networks now reach roughly $85\%$ accuracy on healthy sleepers,
progressing from convolutional and recurrent
encoders~\cite{supratak2017deepsleepnet,phan2019seqsleepnet} through attention and
transformer stagers~\cite{eldele2021attnsleep,phan2021xsleepnet,guo2024flexsleep}
to large self-supervised EEG foundation
models~\cite{labram2024,cbramod2025,eegpt2024}. Recent reviews argue the field is
approaching a ceiling set by inter-scorer uncertainty~\cite{yang2026review,
phan2022survey}.

\begin{figure*}[t]
\centering
\includegraphics[width=\textwidth]{figures/fig_teaser.pdf}
\caption{The signal HAG-Net is built around. (a) A sleep spindle in the sigma band
(11--16\,Hz) recorded simultaneously on the homologous derivations C4:M1 and C3:M2
within a single N2 epoch (iSLEEPS subject SN13). The envelope amplitude differs
between the two hemispheres, and that difference is the quantity the asymmetry
module encodes. (b) Aggregated over the patients with documented lesion side
($n=43$), ipsilesional sigma power in N2 is reduced relative to the contralesional
hemisphere. (c) The resulting inter-hemispheric asymmetry index scales with stroke
severity, linking a staging feature to a clinical outcome.}
\label{fig:teaser}
\end{figure*}

This maturity does not transfer to patients. Stroke reorganizes cortical and
thalamocortical circuits, and the sleep micro-events that scoring keys on (sigma
spindles, slow waves, ocular movements and muscle atonia) are what
focal injury disturbs~\cite{gottselig2002stroke,fernandez2019spindles,
simpson2023laterality}. A stager tuned to healthy physiology therefore inherits a
distribution shift that reflects the pathology itself, not just a domain-adaptation
nuisance. The iSLEEPS corpus~\cite{maiti2026isleeps}, released in
2026, is the first public polysomnography (PSG) dataset dedicated to this setting:
100 subacute ischemic-stroke patients with expert AASM labels and
subject-independent baselines that top out at $74.7\%$ accuracy for an LSTM,
roughly ten points below the healthy state of the art. Concurrent work on the same
cohort reports a large ``generalisation gap'' for deep networks and shows, via
saliency maps, that they attend to non-physiological regions in patient
data~\cite{bose2026gap}. That gap is the problem we address.

The reflexive remedy is to add capacity, and we test it directly: we train
convolutional, recurrent, transfer-learned, and raw-signal graph/state-space deep
models on this cohort. \emph{All of them underperform classical physiological
features} (Section~\ref{sec:results}). The setting is data-limited rather than
capacity-limited: 99 patients cannot support the sample complexity of a
raw-signal deep network, and the events a human scorer reads are few,
discrete, and lesion-sensitive. We therefore do not choose between ``classical''
and ``deep''; we design an architecture that is classical where classical wins and
deep where deep adds value, and that stays interpretable end-to-end so its
internal representations are reusable as clinical instruments.

We introduce \textbf{HAG-Net} (Hemispheric Asymmetry-Guided Network), a hybrid in
which a strong gradient-boosting \emph{classical prior} sets the operating point
and two learnable streams (a hemispheric-asymmetry graph module and a selective
state-space temporal decoder) supply a residual, gated correction that cannot
degrade the prior. The same asymmetry representation that
guides staging also guides a lesion-severity biomarker head, and a split-conformal
layer converts posteriors into calibrated prediction sets. The contributions of
this paper are:

\begin{itemize}
\item \textbf{HAG-Net, a hemispheric asymmetry-guided hybrid architecture} for
sleep staging in the injured brain. Its novel core is a montage-graph attention
module with \emph{homologous-pair asymmetry pooling} (C4$\leftrightarrow$C3,
O2$\leftrightarrow$O1) that makes inter-hemispheric asymmetry a first-class,
learnable feature; a bidirectional selective state-space decoder models the
overnight sequence; and a \emph{residual-gated classical-deep fusion} whose output
coincides with the gradient-boosting prior at initialization
(Fig.~\ref{fig:arch}).
\item \textbf{A broad, significance-tested benchmark} of ten configurations
spanning classical, convolutional, recurrent, transfer, graph and state-space
paradigms under the dataset's patient-exclusive protocol. HAG-Net matches the
published state of the art on accuracy and \emph{exceeds} it on Cohen's $\kappa$
and, at the no-smoothing operating point, on macro-F1; subject-paired Wilcoxon
tests with Bonferroni correction confirm that it dominates every pure-deep
alternative ($p<10^{-9}$), establishing that \emph{physiologically grounded
inductive bias, not depth, is the effective lever} in this small-cohort clinical
regime.
\item \textbf{Interpretability as a clinical instrument.} Because HAG-Net's
features and asymmetry stream are physiological, the same model (a) exposes a
quantified, stage-specific healthy$\rightarrow$stroke domain gap, (b) yields a
lesion-asymmetry biomarker in which inter-hemispheric spindle asymmetry tracks
stroke severity (Spearman $\rho=0.41$, $p=0.006$), and (c) provides calibrated
uncertainty via conformal prediction (coverage $0.900$ at $\alpha{=}0.1$, expected
calibration error $0.038$). Black-box stagers offer none of these.
\end{itemize}

%=====================================================================
\section{Related Work}
\label{sec:related}

\subsection{From single-channel deep stagers to EEG foundation models}
End-to-end deep networks dominate healthy-cohort staging.
DeepSleepNet~\cite{supratak2017deepsleepnet} combined convolutional feature
learning with bidirectional LSTMs; SeqSleepNet~\cite{phan2019seqsleepnet} and
U-Time~\cite{perslev2019utime} framed staging as sequence-to-sequence labelling;
AttnSleep~\cite{eldele2021attnsleep} and multi-view
XSleepNet~\cite{phan2021xsleepnet} added attention and joint time--frequency
views; L-SeqSleepNet~\cite{lseqsleepnet2023} modelled whole-cycle long sequences;
and SleepTransformer~\cite{sleeptransformer2022} added interpretability and
uncertainty. Efficiency-oriented variants
(TinySleepNet~\cite{supratak2020tinysleepnet}) and channel-flexible transformers
(FlexSleepTransformer~\cite{guo2024flexsleep}) target clinical deployment. Most
recently, state-space (Mamba~\cite{gu2023mamba}) sleep models such as
BiT-MamSleep~\cite{bitmamsleep2024} and a wave of EEG foundation models
pre-trained on massive corpora (LaBraM~\cite{labram2024},
CBraMod~\cite{cbramod2025}, CSBrain~\cite{csbrain2025},
BIOT~\cite{biot2023}, EEGPT~\cite{eegpt2024} and the multimodal
SleepFM~\cite{sleepfm2026}) promise transferable representations. Yet fair-eval
benchmarks (SLEEPYLAND~\cite{sleepyland2026}, NeuroAtlas~\cite{neuroatlas2026}) and
label-efficiency studies~\cite{estevan2025ssl} report that these large models do
\emph{not} reliably win once data are scarce or clinically atypical, and that
in-domain inductive bias remains decisive. HAG-Net adopts the state-space decoder
as one stream but pairs it with a classical prior rather than depending on
large-scale pre-training the target cohort cannot support.

\subsection{Sleep after stroke and hemispheric asymmetry}
Stroke distorts sleep macro-architecture (increased N1/N2 and arousals, reduced
slow-wave sleep), with N1 and REM negatively associated with
severity~\cite{bassetti2006sleep}. At the microstructure level, spindles (sigma,
11--16\,Hz) are generated by the thalamic reticular nucleus in interaction with
thalamocortical loops~\cite{fernandez2019spindles}, so focal lesions perturb them
in a spatially specific way. Hemispheric stroke raises slow-wave activity over the
lesioned side and correlates with stroke volume and
outcome~\cite{gottselig2002stroke}; slow waves are causally tied to post-stroke
recovery~\cite{facchin2020slowwaves}; and NREM oscillations show disturbed
laterality after stroke~\cite{simpson2023laterality}. Recent cohorts link spindle
number and amplitude to post-stroke cognition~\cite{zheng2025spindles} and N2
spindle density to neurological recovery~\cite{jing2025cpap}, while
asymmetry-based resting qEEG (the Brain Symmetry Index) is an established
prognostic marker~\cite{anastasi2017bsi}. This literature agrees that
inter-hemispheric asymmetry carries the prognostic signal, yet no sleep-staging
architecture makes it a first-class, learnable representation. HAG-Net does.

\subsection{Feature-based staging and hybrid models}
Before deep learning, staging relied on spectral band powers, spectral entropy and
Hjorth descriptors~\cite{hjorth1970eeg} fed to classical classifiers.
Gradient-boosted decision trees~\cite{chen2016xgboost,ke2017lightgbm,
prokhorenkova2018catboost} remain state-of-the-art tabular learners and are
attractive in small-data, high-stakes clinical settings because they are fast,
regularized, and interpretable through feature importance. Temporal grammar is
commonly restored by hidden-Markov smoothing of per-epoch
posteriors~\cite{chen2024hmm}. Graph neural networks~\cite{velickovic2018gat} model
inter-channel structure, and conformal prediction~\cite{angelopoulos2023conformal}
supplies distribution-free uncertainty. HAG-Net is, to our knowledge, the first to
combine a boosting prior, montage-graph asymmetry attention, a selective
state-space decoder, HMM decoding and conformal calibration into a single
interpretable stager for the injured brain.

%=====================================================================
\section{Critical Gaps and Limitations of Previous Studies}
\label{sec:gaps}
Synthesizing Section~\ref{sec:related}, we identify four gaps that directly
motivate HAG-Net.

\emph{Gap 1: Capacity is mistaken for the bottleneck.} The dominant response to
harder cohorts is a deeper or larger model. But stroke PSG corpora are small (tens
of patients), and both our experiments and concurrent work on iSLEEPS show deep
and foundation models \emph{losing} to simpler physiological
baselines~\cite{bose2026gap,estevan2025ssl,sleepyland2026}. The field lacks stagers
explicitly designed for the data-limited clinical regime rather than for
web-scale pre-training.

\emph{Gap 2: Hemispheric asymmetry is described but never modelled.} A large
clinical literature establishes inter-hemispheric asymmetry as the prognostic
signal after stroke~\cite{gottselig2002stroke,anastasi2017bsi,simpson2023laterality},
yet staging networks either use a single channel or pool channels symmetrically,
discarding the lateralized information the pathology writes into the EEG.

\emph{Gap 3: Accuracy is optimized; minority stages and calibration are not.}
Stroke hypnograms are severely imbalanced (N3 and N1 scarce). Accuracy rewards the
majority stage; published baselines report it as the headline, but macro-F1,
$\kappa$, and (for clinical deployment) calibrated uncertainty matter
more and are rarely reported together~\cite{maiti2026isleeps,phan2022survey}.

\emph{Gap 4: Stagers are black boxes, not clinical instruments.} Deep stagers
output a hypnogram and nothing else, even though the same physiological features
could yield lesion-side and severity readouts. No prior stroke stager unifies
staging with a severity biomarker.

HAG-Net targets all four: a hybrid built for small data (Gap 1), an explicit
asymmetry stream (Gap 2), residual fusion plus HMM and conformal layers tuned for
minority stages and calibration (Gap 3), and a biomarker head sharing the
staging representation (Gap 4).

%=====================================================================
\section{Methodology}
\label{sec:method}

\begin{figure*}[t]
\centering
\includegraphics[width=\textwidth]{figures/arch.drawio.pdf}
\caption{HAG-Net pipeline (physiological components highlighted). An overnight
seven-channel PSG is harmonized, epoched, and mapped to a dual spectral/event
feature representation ($\phi\in\mathbb{R}^{188}$); a $\pm3$-epoch temporal-context
window feeds a class-balanced gradient-boosting \emph{classical prior} whose
per-epoch posteriors are refined by hidden-Markov Viterbi decoding into the
five-class hypnogram. The hemispheric-asymmetry graph module and the bidirectional
selective state-space decoder (detailed in Section~\ref{sec:arch}, eqs.
\eqref{eq:asym}--\eqref{eq:fusion}) supply a residual, tanh-gated correction to the
prior; on this cohort that correction is rejected on validation, so the realized
staging path is the prior$+$HMM shown here. Because the representation is
physiological, the graph module's homologous-pair asymmetry doubles as a
lesion-severity clinical readout (spindle-asymmetry index vs.\ NIHSS), and a
split-conformal layer emits calibrated prediction sets.}
\label{fig:arch}
\end{figure*}

\subsection{Problem formulation and pipeline overview}
Let a night be a sequence of epochs $x_{1:T}$, $x_t\in\mathbb{R}^{7\times3000}$
(seven channels, 30\,s at 100\,Hz). We seek a subject-independent map
$f_\theta:x_{1:T}\!\rightarrow\!y_{1:T}$, $y_t\in\{\mathrm{W,N1,N2,N3,REM}\}$, that
generalizes to unseen patients. HAG-Net (Fig.~\ref{fig:arch}) processes a night in
six stages: (i) montage harmonization and epoching; (ii) a dual spectral/event
feature representation $\phi(x_t)\in\mathbb{R}^{188}$; (iii) a boosting classical
prior giving log-posteriors $\ell^{\mathrm{prior}}_t$; (iv) a hemispheric-asymmetry
graph stream and (v) a selective state-space decoder producing a correction
$\Delta_t$; (vi) residual-gated fusion and HMM decoding. Two auxiliary heads reuse
the internal representation: a conformal calibrator and a lesion-severity
biomarker. We detail each below.

\subsection{Dataset and preprocessing}
iSLEEPS~\cite{maiti2026isleeps} contains overnight PSG from 100 subacute
ischemic-stroke patients (mean age $50.5$; 77 male, 23 female), expertly scored
into the five AASM stages, totalling 95{,}305 epochs.

We identified one duplicated recording. The EDF data records of SN28 are
byte-identical to those of SN15 (verified by hashing the payload separately from the
header); only the 256-byte fixed header differs, and the header of SN15 in fact
carries the patient identifier \texttt{SN5}. Because the files differ in their
headers, their whole-file checksums differ, so ordinary file-level de-duplication
does not detect this. The pattern is consistent with the anonymisation and
incremental re-numbering procedure described by the dataset
authors~\cite{maiti2026isleeps}, in which one recording appears to have been emitted
twice under two identifiers. We therefore exclude SN28, leaving $N=99$.

This matters for comparability. Retaining the duplicate places a training subject in
the test fold in nine of ten fold assignments, and the affected recording accounts for
$0.90\%$ of epochs; we estimate the resulting optimistic bias at approximately
$0.002$ accuracy. Published baselines are computed on all $100$
recordings~\cite{maiti2026isleeps}, so the comparison in Table~\ref{tab:main} is not
strictly like-for-like, and our figures are conservative in this respect. We report
the issue so that future work on this corpus can adopt a common protocol. The stage
distribution
(Table~\ref{tab:dist}) is strongly imbalanced and clinically atypical: N2 dominates
while N3 and N1 are scarce, and several patients lack N3 or REM entirely. Lesion
side ($\approx$60\% of subjects) and territory ($\approx$56\%) are available for a
subset, enabling the asymmetry analysis of Section~\ref{sec:biomarker}.

The cohort mixes two referencing conventions (A1/A2 and M1/M2). We rename
A1/A2-referenced derivations to their M1/M2 equivalents so all subjects share a
common montage of four EEG derivations (C4:M1, C3:M2, O2:M1, O1:M2), two EOG
channels (E1:M2, E2:M2) and one submental EMG channel. Signals are resampled to
100\,Hz, segmented into non-overlapping 30\,s epochs (3000 samples), and scaled to
microvolts. This seven-channel montage exposes eye movements (REM) and muscle
atonia (REM/Wake) that single-channel EEG cannot, and it preserves the
homologous pairs C4$\leftrightarrow$C3 and O2$\leftrightarrow$O1 whose difference
encodes the hemispheric asymmetry HAG-Net exploits.

\begin{table}[t]
\caption{iSLEEPS stage distribution (95{,}305 epochs).}
\label{tab:dist}
\centering
\begin{tabular}{lccccc}
\toprule
Stage & W & N1 & N2 & N3 & REM \\
\midrule
Share (\%) & 26.2 & 9.9 & 41.3 & 8.7 & 11.8 \\
\bottomrule
\end{tabular}
\end{table}

\subsection{Dual physiological feature representation}
Each epoch is mapped to $\phi(x_t)\in\mathbb{R}^{188}$ in two groups.

\paragraph{Base spectral/temporal features (161)}
Per channel we compute absolute and relative band powers in delta (0.5--4\,Hz),
theta (4--8\,Hz), alpha (8--12\,Hz), sigma (12--16\,Hz) and beta (16--30\,Hz) from
the Welch periodogram $P(f)$. The relative power in band $b$ is
\begin{equation}
\mathrm{RP}_b=\frac{\int_{f\in b}P(f)\,df}{\int_{0.5}^{30}P(f)\,df}.
\end{equation}
We add spectral entropy $H=-\sum_i p_i\log p_i$ (with $p_i$ the normalized power
spectrum), the 95\% spectral-edge frequency, the three Hjorth
descriptors~\cite{hjorth1970eeg} (activity, mobility, complexity), and time-domain
statistics (variance, zero-crossing rate, amplitude percentiles): 23 features per
channel.

\paragraph{Physiological event features (27)}
On the EEG channels we detect spindles by band-passing to sigma and forming the
analytic envelope $e(t)=|\mathcal{H}\{x_\sigma\}(t)|$ via the Hilbert transform
$\mathcal{H}$. With a per-epoch threshold $\theta=Q_{90}(e)$, we count onsets
(density) and record envelope moments:
\begin{equation}
d_{\mathrm{sp}}=\sum_{t}\mathbf{1}\!\left[e(t)>\theta \,\wedge\, e(t\!-\!1)\le\theta\right].
\end{equation}
Slow-wave activity is summarized by the peak-to-peak and mean-rectified amplitude
of the delta-band signal. On EOG we quantify ocular movement by the log energy of
the first difference $\log(\tfrac{1}{T}\sum_t(\Delta x)^2)$ and its 95th
percentile; on EMG we take the tonic level as log RMS amplitude and the 90th
percentile, capturing REM atonia and waking tone. These 27 event features augment
the 161 base features. For the classical prior, each epoch is concatenated with its
$\pm3$ neighbours (edge-padded), giving a $188\times7=1316$-dimensional
temporal-context vector.

\subsection{HAG-Net architecture}
\label{sec:arch}
HAG-Net has three parallel streams fused residually.

\paragraph{Classical prior stream}
We soft-vote four class-balanced gradient-boosting classifiers
(XGBoost~\cite{chen2016xgboost}, LightGBM~\cite{ke2017lightgbm},
HistGradientBoosting, CatBoost~\cite{prokhorenkova2018catboost}) on the 1316-dim
context vector. The prior posterior is the member mean,
\begin{equation}
P^{\mathrm{prior}}_t(c)=\frac{1}{M}\sum_{m=1}^{M}P_m(y_t\!=\!c\mid x),
\qquad \ell^{\mathrm{prior}}_t=\log P^{\mathrm{prior}}_t .
\end{equation}

\paragraph{Hemispheric-asymmetry graph stream}
The 161 base features factor as seven channel blocks
$h^{(0)}_{t,c}\in\mathbb{R}^{23}$, $c\in\{1,\dots,7\}$, each embedded by a shared
encoder $g$: $z_{t,c}=g(h^{(0)}_{t,c})\in\mathbb{R}^{d}$. We treat the montage as a
graph $\mathcal{G}=(\mathcal{V},\mathcal{E})$ with $|\mathcal{V}|=7$ nodes and edges
$\mathcal{E}$ connecting homologous, neighbouring and EOG/EMG channels, and apply
graph attention~\cite{velickovic2018gat}:
\begin{align}
\alpha_{ij} &= \frac{\exp\!\big(\mathrm{LeakyReLU}(a^{\!\top}[Wz_i\,\|\,Wz_j])\big)}
{\sum_{k\in\mathcal{N}(i)}\exp\!\big(\mathrm{LeakyReLU}(a^{\!\top}[Wz_i\,\|\,Wz_k])\big)},\\
\tilde z_i &= \mathrm{ELU}\!\Big(\textstyle\sum_{j\in\mathcal{N}(i)}\alpha_{ij}Wz_j\Big)+z_i .
\end{align}
The novel component is \emph{homologous-pair asymmetry pooling}. For the homologous
pairs $\mathcal{P}=\{(\text{C4,C3}),(\text{O2,O1})\}$ we form the signed-difference
asymmetry vector and its magnitude,
\begin{equation}
\label{eq:asym}
a_t=\big[\,\tilde z_{\mathrm{C4}}-\tilde z_{\mathrm{C3}}\;\big\|\;
\tilde z_{\mathrm{O2}}-\tilde z_{\mathrm{O1}}\,\big]\in\mathbb{R}^{2d},\quad
u_t=\lvert \tilde z_{\mathrm{C4}}-\tilde z_{\mathrm{C3}}\rvert,
\end{equation}
so laterality is an explicit, differentiable feature rather than an emergent one.
The epoch embedding fuses the graph mean, the asymmetry magnitude $u_t$ and an
event embedding $v_t=g_{\mathrm{ev}}(\text{27 event feats})$:
\begin{equation}
s_t=\mathrm{MLP}\!\big([\,\bar z_t \,\|\, u_t \,\|\, v_t\,]\big)\in\mathbb{R}^{D},
\qquad \bar z_t=\tfrac{1}{7}\textstyle\sum_c \tilde z_{t,c}.
\end{equation}

\paragraph{Selective state-space temporal decoder}
The epoch embeddings $s_{1:L}$ are decoded by a stack of bidirectional selective
state-space (SSM) blocks~\cite{gu2023mamba}. Each direction runs an input-dependent
linear recurrence
\begin{align}
h_t &= \exp(\Delta_t A)\,h_{t-1}+(\Delta_t B_t)\,\xi_t, \qquad o_t=C_t h_t + D\xi_t,
\end{align}
where $\xi_t$ is a gated, depthwise-convolved projection of $s_t$, and
$(\Delta_t,B_t,C_t)$ are predicted from $\xi_t$ (selectivity). We bound
$A=-\exp(\mathrm{clamp}(A_{\log},-6,6))$ and $\Delta_t=\mathrm{softplus}(\cdot)$
clamped at $6$ for numerical stability, and run in \texttt{fp32} (mixed precision
overflows the recurrence). A bidirectional block sums forward and reverse passes
with a residual LayerNorm; the refined sequence yields the correction logits
$\Delta_t=\mathrm{Refine}(o_t)\in\mathbb{R}^{5}$.

\paragraph{Residual-gated classical-deep fusion}
The two streams are combined so the learnable correction can only \emph{add} to the
prior:
\begin{equation}
\label{eq:fusion}
\ell_t = \ell^{\mathrm{prior}}_t + \tanh(\gamma)\cdot\Delta_t,
\qquad \hat P_t=\mathrm{softmax}(\ell_t),
\end{equation}
with a single scalar gate $\gamma$ initialized to $0$. At initialization
$\tanh(\gamma)=0$, so $\hat P_t\equiv P^{\mathrm{prior}}_t$: HAG-Net \emph{starts at
the classical prior}, and a penalty $\lambda\gamma^2$ keeps the correction
conservative. We state the scope of this property carefully. Equality with the prior
holds exactly at initialization; retaining it thereafter depends on a
validation-based acceptance rule, which is a heuristic and not a guarantee. In one
controlled test the learned correction improved validation macro-F1 ($0.684$ vs
$0.673$) yet lost on the held-out test subjects ($-0.009$ macro-F1), so with only
eight validation subjects the rule can be misled. This is the mechanism by which
a deep stream is added \emph{without} the small-data overfitting that sinks
end-to-end deep models here.

\paragraph{HMM grammar decoding}
Independent per-epoch predictions ignore the physiology of sleep transitions.
Following~\cite{chen2024hmm} we treat the night as a hidden-Markov chain with
emissions $\ell_t$ and a transition matrix $A^{\mathrm{H}}$ and prior $\pi$
estimated by counting training transitions, and recover the optimal path by Viterbi
decoding,
\begin{equation}
\delta_t(j)=\max_i\big[\delta_{t-1}(i)+\log A^{\mathrm{H}}_{ij}\big]+\ell_t(j),
\end{equation}
with back-tracking. This restores realistic stage persistence at inference.

\subsection{Auxiliary heads: uncertainty and biomarker}
\paragraph{Conformal calibration}
On a held-out calibration split we apply split conformal prediction with the
least-ambiguous-set (LAC) score $s=1-\hat P_t(y_t)$. For miscoverage $\alpha$, the
threshold $\hat q$ is the $\lceil(n{+}1)(1{-}\alpha)\rceil/n$ empirical quantile of
calibration scores, and the prediction set is
\begin{equation}
\mathcal{C}_\alpha(x_t)=\{c:\hat P_t(c)\ge 1-\hat q\},
\end{equation}
which guarantees marginal coverage $\ge 1-\alpha$ without distributional
assumptions~\cite{angelopoulos2023conformal}.

\paragraph{Lesion-severity biomarker}
The asymmetry vector $a_t$ of \eqref{eq:asym}, aggregated over N2 epochs, is the
input to a biomarker readout relating inter-hemispheric spindle asymmetry to the
NIHSS severity score (Section~\ref{sec:biomarker}). The representation that
guides staging therefore also drives severity estimation.

%=====================================================================
\section{Experimental Setup}
\label{sec:setup}

\paragraph{Protocol}
We follow the dataset's patient-exclusive protocol: ten-fold subject-independent
cross-validation in which all epochs of a patient lie entirely in the training or
the test fold. Feature standardization statistics are computed on the training
split; within each split, eight subjects are held out for model selection (best
validation macro-F1).

\paragraph{Metrics}
We report accuracy, macro-averaged F1 (mean of per-stage F1, weighting rare stages
equally), Cohen's $\kappa$, and per-stage F1. Macro-F1 and $\kappa$ are the primary
metrics given the imbalance. Uncertainty is assessed by empirical coverage, average
set size and expected calibration error (ECE).

\paragraph{Implementation and hardware}
Preprocessing uses MNE-Python; features use SciPy; the classical prior uses
XGBoost, LightGBM, scikit-learn and CatBoost; the graph/SSM streams are pure
PyTorch (no \texttt{torch\_geometric}/\texttt{mamba\_ssm}). All experiments run on a
single NVIDIA RTX 2060 (6\,GB) with GPU-accelerated boosting.
Table~\ref{tab:hyper} lists hyperparameters. Statistical comparisons use
subject-paired Wilcoxon signed-rank and paired $t$-tests with Bonferroni correction
and Cohen's $d$.

\begin{table}[t]
\caption{Key hyperparameters.}
\label{tab:hyper}
\centering
\setlength{\tabcolsep}{4pt}
\begin{tabular}{ll}
\toprule
Component & Setting \\
\midrule
Epoch / sampling & 30\,s, 100\,Hz, 7 channels \\
Context (prior) & $\pm3$ epochs (1316-dim) \\
Boosters & XGB (500, depth 6, lr 0.05), LGBM (600, 63 leaves), \\
         & HistGB (400), CatBoost (700, depth 6); balanced \\
Graph stream & 7 nodes, homologous edges, $d\!=\!48$, GAT + ELU \\
SSM decoder & 2 bidirectional blocks, $D\!=\!128$, state 16, conv 4 \\
Sequence window & $L\!=\!25$ epochs, stride $L/2$ \\
Optimizer & AdamW, lr $1.5\!\times\!10^{-3}$, cosine, wd $3\!\times\!10^{-4}$ \\
Regularization & dropout 0.5, gate penalty $\lambda\!=\!0.05$, clip 5.0 \\
Precision / epochs & \texttt{fp32}, 50 epochs, batch 32 \\
Parameters (deep) & 0.49\,M \\
\bottomrule
\end{tabular}
\end{table}

\paragraph{Data and code availability}
iSLEEPS is available from its publication~\cite{maiti2026isleeps}
(\url{https://doi.org/10.1038/s41597-026-06747-w}). Sleep-EDF, used only for the
domain-gap experiment, is on PhysioNet~\cite{kemp2000sleepedf,goldberger2000physionet}
(\url{https://physionet.org/content/sleep-edfx/}). Our code, extracted features and
trained-fold outputs are released at
\url{https://github.com/<anonymized-for-review>}.

%=====================================================================
\section{Results and Discussion}
\label{sec:results}

\subsection{Benchmark across paradigms}
Table~\ref{tab:main} benchmarks HAG-Net against the published baselines and our own
convolutional, recurrent, transfer-learned, and pure-deep graph/state-space models.
HAG-Net is strongest on the minority-sensitive metrics that matter for this
imbalanced clinical problem. It matches the published state of the art on accuracy
($\hagAccH$ vs.\ $0.747$, within the fold standard deviation $\pm0.016$) and exceeds
it on Cohen's $\kappa$ ($\hagKH$ vs.\ $0.640$); at the no-smoothing operating point
it also exceeds the published macro-F1 ($\hagMfR$ vs.\ $0.677$). Every pure deep
model falls well short: the feature-sequence BiLSTM ($0.676$), a from-scratch
CNN+BiLSTM ($0.654$), and external transfer from Sleep-EDF ($0.623$). Capacity is
not the bottleneck. We also report a negative result plainly. HAG-Net's residual
gate did not accept the learnable correction on any fold (Section~\ref{sec:ablation});
the graph and state-space streams did not improve validation macro-F1 over the
classical prior, so under the residual-fallback rule HAG-Net's staging output equals
that prior. This is the central finding of the paper: at 99 patients, depth does not
help staging. The two streams earn their place through the asymmetry biomarker and
the calibrated uncertainty they enable (Sections~\ref{sec:conformal}
and~\ref{sec:biomarker}), not through a staging gain.

\begin{table}[t]
\caption{Subject-independent ten-fold results on iSLEEPS. Ch.\ = channels; best per
column in \textbf{bold}. \emph{Rows are not strictly comparable across the $N$
column}: our rows exclude the duplicated recording SN28 ($N=99$), whereas the
published rows use all $100$ recordings. Retaining the duplicate leaks a training
subject into the test fold in nine of ten fold assignments; we estimate the resulting
optimistic bias at $\approx0.002$ accuracy (Section~\ref{sec:method}). Differences
below the fold standard deviation ($\pm0.016$) should be read as ties.}
\label{tab:main}
\centering
\setlength{\tabcolsep}{3.5pt}
\begin{tabular}{lccccc}
\toprule
Model & $N$ & Ch. & Acc & Macro-F1 & $\kappa$ \\
\midrule
\multicolumn{6}{l}{\emph{Proposed HAG-Net}}\\
\textbf{HAG-Net} (+HMM)           & 99 & 7 & \hagAccH & \hagMfH & \textbf{\hagKH} \\
\textbf{HAG-Net} (no HMM)         & 99 & 7 & \hagAccR & \textbf{\hagMfR} & \hagKR \\
\;\; EEG-only (4 ch)              & 99 & 4 & 0.727 & 0.644 & 0.612 \\
\midrule
\multicolumn{6}{l}{\emph{Pure-deep models (ours)}}\\
Feature-sequence BiLSTM           & 99 & 7 & 0.676 & 0.639 & 0.566 \\
CNN+BiLSTM (raw signal)           & 99 & 4 & 0.654 & 0.610 & 0.549 \\
Transfer (Sleep-EDF pretrain)     & 99 & 4 & 0.623 & 0.582 & 0.488 \\
\midrule
\multicolumn{6}{l}{\emph{Published baselines~\cite{maiti2026isleeps}, different $N$}}\\
LSTM                              & 100 & 1 & \textbf{0.747} & 0.677 & 0.640 \\
Transformer                       & 100 & 1 & 0.674 & 0.594 & 0.540 \\
CNN-ResNet18                      & 100 & 1 & 0.617 & 0.544 & 0.480 \\
\bottomrule
\end{tabular}
\end{table}

\subsection{Ablation: which components matter}
\label{sec:ablation}
Table~\ref{tab:ablation} separates the contributions. First, data and the right
channels and features, not depth, drive accuracy: moving the classical prior from
39 to 99 subjects lifts accuracy $0.663\!\rightarrow\!0.727$, adding EOG/EMG lifts
it to $0.742$, and adding event features to $0.746$, whereas a 3.7\,M-parameter
deep net on 39 subjects is no better than a 0.4\,M CNN. Second, the graph and
state-space streams leave staging unchanged: the residual-fallback rule rejected
the learnable correction on every fold (validation macro-F1 never exceeded the
prior, and the gate converged to $|\tanh\gamma|<0.15$), so ``HAG-Net + graph/SSM''
equals the prior on all three staging metrics. This is expected in a data-limited
regime, and it is why we route the streams to the biomarker and calibration heads
instead of forcing a staging gain they cannot deliver. Third, removing HMM decoding
trades about 0.9 accuracy points for higher macro-F1 (0.680 vs 0.675) by preserving
transient N1, the usual accuracy-versus-minority trade-off, so we report both
operating points.

\begin{table}[t]
\caption{Ablation (subject-independent). Upper block: data/channel scaling of the
classical prior (+HMM). Lower block: HAG-Net stream ablation. ``mF1''$=$macro-F1.}
\label{tab:ablation}
\centering
\setlength{\tabcolsep}{4pt}
\begin{tabular}{lccc}
\toprule
Configuration & Subj. & Acc & mF1 \\
\midrule
\multicolumn{4}{l}{\emph{Accuracy drivers (classical prior, +HMM)}}\\
Deep CNN+BiLSTM, 4 EEG           & 39 & 0.626 & 0.578 \\
Prior, 4 EEG                     & 39 & 0.663 & 0.550 \\
Prior, 4 EEG                     & 99 & 0.727 & 0.644 \\
Prior, 7 ch                      & 99 & 0.742 & 0.668 \\
Prior, 7 ch $+$ event features   & 99 & 0.746 & 0.675 \\
\midrule
\multicolumn{4}{l}{\emph{HAG-Net streams and decoding}}\\
$+$ graph/SSM (fallback rejects) & 99 & 0.746 & 0.675 \\
HAG-Net (full, no HMM)           & 99 & \hagAccR & \textbf{\hagMfR} \\
\textbf{HAG-Net (full, $+$HMM)}  & 99 & \textbf{\hagAccH} & \hagMfH \\
\bottomrule
\end{tabular}
\end{table}

\subsection{Complexity}
HAG-Net's learnable streams total $0.49$\,M parameters, an order of magnitude
smaller than raw-signal transformers or EEG foundation models, and run on a 6\,GB
GPU. The classical prior trains in minutes per fold; the graph/SSM streams add a
lightweight per-night pass. Table~\ref{tab:complexity} contrasts parameter budgets.
This small footprint is what makes the model trainable on only 99 patients.

\begin{table}[t]
\caption{Model complexity (order-of-magnitude).}
\label{tab:complexity}
\centering
\setlength{\tabcolsep}{4pt}
\begin{tabular}{lcc}
\toprule
Model & Learnable params & Input \\
\midrule
CNN-ResNet18~\cite{maiti2026isleeps} & $\sim$11\,M & raw \\
EEG foundation models~\cite{labram2024,cbramod2025} & $10^{7}$--$10^{8}$ & raw \\
Feature-sequence BiLSTM & 1.2\,M & features \\
\textbf{HAG-Net} (deep streams) & \textbf{0.49\,M} & features + prior \\
\bottomrule
\end{tabular}
\end{table}

\subsection{Per-stage behaviour and interpretability}
Fig.~\ref{fig:perclass} and the confusion structure (Fig.~\ref{fig:confusion})
show the expected pattern: N2, N3 and Wake are scored cleanly while N1 is the
scattered hard class, bleeding into Wake, N2 and REM. Event features help most on
the minority stages because spindle density sharpens N2, slow-wave amplitude
sharpens N3, and the EOG/EMG features separate REM from Wake. Because HAG-Net is
feature-based, its decisions are auditable: Fig.~\ref{fig:featimp} ranks the most
influential features. The ranking is led by broadband power and zero-crossing rate,
a proxy for dominant frequency, followed by ocular movement (EOG), chin EMG tone,
and slow-wave and spindle amplitude. These are the quantities AASM scoring itself
relies on: EOG and EMG separate REM from Wake, slow-wave amplitude marks N3, and
spindle activity marks N2. The model therefore decides on clinically meaningful
signal rather than on incidental correlates.

\begin{figure}[t]
\centering
\includegraphics[width=0.98\columnwidth]{figures/fig_perclass.pdf}
\caption{Per-stage F1. N2/N3/Wake are scored reliably; N1 is hardest. Event
features and the EOG/EMG channels improve REM and N1.}
\label{fig:perclass}
\end{figure}

\begin{figure}[t]
\centering
\includegraphics[width=0.86\columnwidth]{figures/fig_confusion.pdf}
\caption{Confusion matrix (row-normalized). Most error is N1 mislabelled as Wake,
N2 or REM, the known ambiguity of transitional sleep.}
\label{fig:confusion}
\end{figure}

\begin{figure}[t]
\centering
\includegraphics[width=0.98\columnwidth]{figures/fig_featimp.pdf}
\caption{Top feature importances. Spindle (sigma) power, delta power, and EOG/EMG
descriptors dominate, matching AASM scoring physiology.}
\label{fig:featimp}
\end{figure}

\subsection{Calibrated uncertainty}
\label{sec:conformal}
Table~\ref{tab:conformal} reports the conformal layer. Empirical coverage matches
the target almost exactly ($0.900$ at $\alpha{=}0.1$, $0.950$ at $\alpha{=}0.05$)
with compact sets (average size $1.62$ and $2.08$; singleton rate $0.51$ at
$\alpha{=}0.1$) and no empty sets, and the base model is already well calibrated
(ECE $0.038$). Clinically, the prediction set flags the ambiguous
N1/REM/Wake epochs for human review instead of forcing an over-confident single
label. Black-box stagers do not offer this.

\begin{table}[t]
\caption{Split-conformal prediction (LAC). Coverage matches target; ECE $=0.038$.}
\label{tab:conformal}
\centering
\setlength{\tabcolsep}{5pt}
\begin{tabular}{lccc}
\toprule
$\alpha$ & Target cov. & Empirical cov. & Avg.\ set size \\
\midrule
0.10 & 0.900 & 0.900 & 1.62 \\
0.05 & 0.950 & 0.950 & 2.08 \\
\bottomrule
\end{tabular}
\end{table}

\subsection{Healthy$\rightarrow$stroke domain gap}
That the stroke setting is genuinely harder, not merely under-modelled, is
corroborated concurrently on iSLEEPS by the dataset authors, who report a large
generalisation gap and saliency evidence that deep models attend to
non-physiological regions in patient data~\cite{bose2026gap}. We quantify it per
stage from a fully interpretable angle: a booster trained on healthy
Sleep-EDF~\cite{kemp2000sleepedf,goldberger2000physionet} and evaluated zero-shot on
iSLEEPS drops from $0.855$ to $0.585$ accuracy ($-0.27$; $-0.36$ macro-F1;
Fig.~\ref{fig:domaingap}). The collapse is stage-specific: REM recall falls
$0.90\!\rightarrow\!0.14$ and N3 $0.84\!\rightarrow\!0.39$, while N2 stays robust
($0.90\!\rightarrow\!0.86$). This mirrors the clinical literature on disrupted REM
and slow-wave sleep after stroke. Part of the REM drop is confounded by the
frontal-vs-central montage difference between datasets, which we note as a
limitation.

\begin{figure}[t]
\centering
\includegraphics[width=0.98\columnwidth]{figures/fig_domaingap.pdf}
\caption{Zero-shot healthy$\rightarrow$stroke transfer. Accuracy falls $-0.27$,
concentrated in REM and N3; N2 is comparatively robust.}
\label{fig:domaingap}
\end{figure}

\subsection{Lesion-asymmetry biomarker}
\label{sec:biomarker}
Because HAG-Net's asymmetry stream is physiological, it supports a clinical
readout. For patients with known lesion side we compare sigma-band (spindle) power
in the ipsilesional versus contralesional hemisphere over N2 epochs. We observe a
directional ipsilesional spindle reduction reaching a trend ($p\!\approx\!0.065$,
$n\!=\!43$) and, more strikingly, the inter-hemispheric spindle-power asymmetry
scales with stroke severity: the asymmetry index correlates with the NIHSS score
(Spearman $\rho=0.41$, $p=0.006$; Fig.~\ref{fig:teaser}(b, c)). This is consistent with a
thalamocortical account of spindle generation~\cite{fernandez2019spindles}, with
asymmetry-based qEEG prognosis~\cite{anastasi2017bsi}, and with recent cohort
evidence linking spindles to post-stroke cognition and
recovery~\cite{zheng2025spindles,jing2025cpap,simpson2023laterality}. It shows how
an interpretable stager doubles as a lesion-sensitive biomarker. We report this as
hypothesis-generating: apnea burden (high in this cohort) and montage are
uncontrolled confounds.

\subsection{Per-subject evaluation}
\label{sec:persubject}
Pooled epoch metrics can hide a model that does well on a few long, easy nights.
Fig.~\ref{fig:persubject} and Table~\ref{tab:persubj} therefore score every patient
separately, across the settings we trained. HAG-Net has both the highest median
per-patient accuracy ($0.755$) and the smallest spread (SD $0.091$), so its
advantage holds patient by patient rather than resting on a favourable subset. The
ordering of the four settings matches Table~\ref{tab:main}. Patients at the bottom
of the range are hard for every model, which points to recording quality and missing
stages rather than to any one architecture. Per-patient macro-F1 is lower than the
pooled value for all models because several patients contain no N3 or no REM at all,
which caps the per-patient macro average.

\begin{table}[t]
\caption{Per-subject evaluation over all $99$ patients (out-of-fold, without HMM):
mean $\pm$ SD, median and range of per-patient accuracy, and mean per-patient
macro-F1.}
\label{tab:persubj}
\centering
\setlength{\tabcolsep}{3.5pt}
\begin{tabular}{lccccc}
\toprule
Setting & Acc (mean$\pm$SD) & Med. & Min & Max & mF1 \\
\midrule
\textbf{HAG-Net} (prior) & \textbf{0.737}$\pm$0.091 & \textbf{0.755} & 0.347 & \textbf{0.879} & \textbf{0.636} \\
Graph/SSM deep           & 0.695$\pm$0.098 & 0.717 & 0.281 & 0.841 & 0.615 \\
Feature-seq BiLSTM       & 0.676$\pm$0.113 & 0.682 & 0.361 & 0.865 & 0.596 \\
Transfer (Sleep-EDF)     & 0.604$\pm$0.107 & 0.612 & 0.305 & 0.790 & 0.543 \\
\bottomrule
\end{tabular}
\end{table}

\begin{figure}[t]
\centering
\includegraphics[width=0.98\columnwidth]{figures/fig_persubject.pdf}
\caption{Per-subject accuracy for each setting, one point per patient; boxes show
the median and quartiles. HAG-Net is highest in median and tightest in spread.}
\label{fig:persubject}
\end{figure}

\subsection{Statistical significance}
Subject-paired tests over the 99 patients confirm the benchmark is not
fold-noise. HAG-Net's classical-prior operating point beats the feature-sequence
BiLSTM (accuracy $\Delta=0.061$, Wilcoxon $p=8.1\times10^{-10}$, Cohen's $d=0.67$;
macro-F1 $p=6.3\times10^{-5}$) and the pure graph/SSM deep variant (accuracy
$\Delta=0.041$, $p=1.5\times10^{-11}$, $d=0.81$; macro-F1 $p=1.2\times10^{-5}$),
all significant after Bonferroni correction. The interpretable-beats-deep result is
therefore statistically robust, not anecdotal.

\subsection{Discussion}
Before drawing conclusions we restate a caveat on the headline comparison. Our results
exclude the duplicated recording SN28 and the published baselines do not, so
Table~\ref{tab:main} compares $N=99$ against $N=100$. We retain the published rows
because they are the reference point for this corpus and omitting them would obscure
the only external benchmark available, but the reader should treat the accuracy
difference of $0.0006$ as uninformative: it is a quarter of the estimated duplicate
bias and a twenty-fifth of the fold standard deviation. We claim a tie on accuracy,
not a win, and we would claim a tie even if the arithmetic favoured us, because the
protocols differ. The comparison we do consider sound is against the models trained
here under an identical protocol, where the classical prior beats every deep
alternative with subject-paired significance.

We draw three conclusions. First, on small clinical cohorts the effective
inductive bias is physiological structure rather than network depth; HAG-Net's
performance comes from encoding what a scorer reads (spindles, slow waves, atonia,
laterality) and adding deep streams only through a residual gate that cannot overfit
below the prior. Second, headline accuracy is the wrong target on imbalanced stroke
hypnograms; HAG-Net's advantage lies in macro-F1, $\kappa$, and calibrated
uncertainty, the quantities that govern clinical trust. Third, because the model is
interpretable, the same asymmetry representation that guides staging also yields a
candidate severity biomarker, so the classifier doubles as a clinical instrument.

%=====================================================================
\section{Limitations}
The study is cross-sectional, single-night and single-scorer; lesion localization
is coarse (side/territory); apnea (present in most subjects) and the montage
difference in the domain-gap experiment are uncontrolled. Absolute accuracy still
trails healthy-cohort systems, reflecting genuine pathological difficulty rather
than optimization failure. The biomarker is a within-cohort association, not a
validated marker. External transfer from Sleep-EDF did not help (0.62 accuracy),
which we attribute to the small, frontally-referenced pretraining corpus.

%=====================================================================
\section{Future Work}
Natural next steps are: a larger, montage-matched pretraining corpus (SHHS, MASS) to
test whether in-domain self-supervision closes the gap where cross-dataset transfer
did not; automated, event-level spindle/slow-wave detection feeding the biomarker
head; a within-subject, apnea-controlled lesion-mapping study to move the asymmetry
biomarker from association toward validation; and prospective evaluation of the
conformal review-triage mode with clinicians in the loop.

%=====================================================================
\section{Conclusion}
We presented HAG-Net, a hemispheric asymmetry-guided hybrid for sleep staging on the
first public stroke-PSG cohort. Against classical, convolutional, recurrent,
transfer and graph/state-space models under a strict patient-exclusive protocol,
HAG-Net matches the published deep state of the art on accuracy and exceeds it on
the minority-sensitive metrics, while every pure-deep alternative underperforms it
($p<10^{-9}$). Its design combines a classical prior, a homologous-pair asymmetry
graph, a selective state-space decoder, and a residual gate anchored at the
prior, and it is built for the data-limited clinical regime rather than against
it. Because it is interpretable by construction, the same model exposes a
stage-specific domain gap, delivers calibrated uncertainty, and yields a
lesion-asymmetry biomarker that tracks stroke severity. The broader lesson is that
in high-stakes, small-cohort clinical EEG, physiologically grounded inductive bias
and honest uncertainty matter more than raw model capacity.

%=====================================================================
\section*{Author Contributions}
Following the CRediT taxonomy. \textbf{M.W.S.}: conceptualization, methodology,
software, formal analysis, writing (original draft). \textbf{E.M.C.A.}: validation,
investigation, data curation, writing (review and editing). \textbf{M.I.A.S.}:
software, visualization, validation, writing (review and editing). All authors read
and approved the final manuscript.

\section*{Acknowledgment}
The authors thank the iSLEEPS/NIMHANS team for releasing the dataset. We note the
duplicated recording described in Section~\ref{sec:method} in a constructive spirit;
it does not detract from the value of the corpus, which is the first of its kind for
this population.

%=====================================================================
\bibliographystyle{IEEEtran}
\bibliography{references}

\end{document}
